\chapter{Przykłady elementów pracy dyplomowej}
\label{cha:przyklady}

\section{Liczba}

Pakiet \texttt{siunitx} zadba o to, by liczba została poprawnie sformatowana: \\
\begin{center}
	\num{1234567890.0987654321}
\end{center}


\section{Rysunek}

Pakiet \texttt{subcaption} pozwala na umieszczanie w podpisie rysunku odnośników do ,,podilustracji'': \\

\begin{figure}[h]
	\centering
	\begin{subfigure}{0.35\textwidth}
		\centering
		\framebox[2.0\width]{A}
		\subcaption{\label{subfigure_a}}
	\end{subfigure}
	\begin{subfigure}{0.35\textwidth}
		\centering
		\framebox[2.0\width]{B}
		\subcaption{\label{subfigure_b}}
	\end{subfigure}

	\caption{\label{fig:subcaption_example}Przykład użycia \texttt{\textbackslash subcaption}: \protect\subref{subfigure_a} litera A, \protect\subref{subfigure_b} litera B.}
\end{figure}

Opis rysunku umieszczamy pod rysunkiem. Do każdego zamieszczonego rysunku należy odwołać się w tekście co najmniej raz. Rysunek powinien zostać zamieszczony w kodzie dokumentu bezpośrednio po akapicie, w którym po raz pierwszy wystąpiło do niego odwołanie.

Aby rysunki nie ,,uciekały'' do kolejnych rozdziałów, na końcu rozdziału (lub sekcji) można wstawić polecenie \texttt{\textbackslash FloatBarrier}.

\section{Tabela}

Pakiet \texttt{threeparttable} umożliwia dodanie do tabeli adnotacji: \\

\begin{table}[h]
	\centering

	\begin{threeparttable}
		\caption{Przykład tabeli}
		\label{tab:table_example}

		\begin{tabularx}{0.6\textwidth}{C{1}}
			\toprule
			\thead{Nagłówek\tnote{a}} \\
			\midrule
			Tekst 1 \\
			Tekst 2 \\
			\bottomrule
		\end{tabularx}

		\begin{tablenotes}
			\footnotesize
			\item[a] Jakiś komentarz\textellipsis
		\end{tablenotes}

	\end{threeparttable}
\end{table}

Opis tabeli umieszczamy \textbf{nad} tabelą. Podobnie jak w przypadku rysunków, do każdej zamieszczonej tabeli należy odwołać się w tekście przynajmniej raz.

\section{Wzory matematyczne}

Czasem zachodzi potrzeba wytłumaczenia znaczenia symboli użytych w równaniu. Można to zrobić z użyciem zdefiniowanego na potrzeby niniejszej klasy środowiska \texttt{eqwhere}.
%
\begin{equation}
E = mc^2
\end{equation}
gdzie % ważne, żeby nie zostawić tu pustej linii
\begin{eqwhere}[2cm]
	\item[$m$] masa
	\item[$c$] prędkość światła w próżni
\end{eqwhere}

Odległość półpauzy od lewego marginesu należy dobrać pod kątem najdłuższego symbolu (bądź listy symboli) poprzez odpowiednie ustawienie parametru tego środowiska (domyślnie: 2 cm).

W przypadku przekształceń czy wyprowadzeń na kilka linii, możemy zastosować środowisko \texttt{eqnarray}, które pozwoli wyrównać równania, np. do znaku równości.
%
\begin{eqnarray}
    a - b &=& c\label{eq:roznica}\\
    a &=& b + c\label{eq:wynik}
\end{eqnarray}

Do równań można się odwołać poleceniem \texttt{\textbackslash eqref}. Podstawiając \eqref{eq:roznica} do \eqref{eq:wynik} otrzymamy tożsamość.

Środowisko matematyczne niektóre litery traktuje jako szczególne symbole i dodaje do nich odstęp. Należy do nich m.in. litera ,,N'', dlatego też w zapisie $NF$ (\texttt{\$NF\$}) widoczny jest dodatkowy odstęp pomiędzy $N$ a $F$. Aby tego uniknąć, należy wszystkie litery danego oznaczenia zamknąć w odpowiednim środowisku, np. \texttt{mathit}, \texttt{mathrm} lub \texttt{text}. Jesli chcemy dodatkowo usunąć odstęp, możemy skorzystać z polecenia \texttt{\textbackslash !}. Wyniki prezentuje tabela~\ref{tab:math_acronym}.

\begin{table}
    \centering
    \caption{Zapis symbolu złożonego z kilku liter w środowisku matematycznym}
    \label{tab:math_acronym}
	\begin{tabular}{ll}
        \toprule
        Kod & Efekt \\
        \midrule
		\texttt{\$NF\$} & $NF$ \\
		\texttt{\$\textbackslash mathit\{NF\}\$} & $\mathit{NF}$ \\
		\texttt{\$\textbackslash mathrm\{NF\}\$} & $\mathrm{NF}$ \\
		\texttt{\$\textbackslash text\{NF\}\$}   & $\text{NF}$ \\
        \texttt{\$f\_0\$} & $f_0$ \\
		\texttt{\$f\textbackslash!\_0\$} & $f\!_0$ \\
        \bottomrule
	\end{tabular}
\end{table}

Jeśli symbolu używamy często, dla uproszczenia można zdefiniować własne polecenie w preambule dokumentu, np. \texttt{\textbackslash newcommand\{\textbackslash NF\}\{\textbackslash mathit\{NF\}\}}.
